\section{Document Structure}\label{sec:document-structure}

\subsection{Sections and Subsections}\label{subsec:sections-subsections}
LaTeX provides a hierarchical structure for organizing documents through sections, subsections, and subsubsections.

\subsubsection{Creating Sections}\label{subsubsec:creating-sections}
Sections are created using the \verb|\section{}| command.
The text inside the braces becomes the title of the section.

\subsubsection{Section Hierarchy}\label{subsubsec:section-hierarchy}
LaTeX supports up to 7 levels of depth:
\begin{itemize}
    \item \verb|\part{}| - Only used in large documents
    \item \verb|\chapter{}| - Only available in book and report classes
    \item \verb|\section{}| - Main divisions of a document
    \item \verb|\subsection{}| - Subdivisions of a section
    \item \verb|\subsubsection{}| - Further subdivisions
    \item \verb|\paragraph{}| - Rarely used
    \item \verb|\subparagraph{}| - Rarely used
\end{itemize}

\subsection{Labels and Cross-References}\label{subsec:labels-cross-refs}
Labels allow you to reference sections, figures, tables, and more throughout your document.

\subsubsection{Creating Labels}\label{subsubsec:creating-labels}
Add a label using \verb|\label{sec:my-label}| after a section heading or inside a figure/table environment.

\subsubsection{Referencing Labels}\label{subsubsec:referencing-labels}
Reference a label using \verb|\ref{sec:my-label}| to display its number, or \verb|\pageref{sec:my-label}| to display its page number.

Example: This section is Section~\ref{sec:document-structure} on page~\pageref{sec:document-structure}.

\subsection{Table of Contents}\label{subsec:table-of-contents}
The table of contents is generated automatically from section headings using the \verb|\tableofcontents| command (as seen at the beginning of this document).

\subsubsection{Customizing TOC Depth}\label{subsubsec:toc-depth}
Control which levels appear in the TOC with:
\begin{verbatim}
\setcounter{tocdepth}{2}  % Show sections and subsections
\end{verbatim}